% Source: Weinberg GR, physical PDF pp. 47--50
%         (printed pp. 20--23).
% Coverage: Chapter 1 bibliography and references 1--40.
% Figures/tables/footnotes: none.
% Status: source-reviewed and compile-clean.
% Uncertainties: none.

\chapterbackmatter{Bibliography}

Not being an historian, I have been content to base this chapter on secondary
sources, aside from the works of Newton, Mach, Maxwell, Newcomb, and Einstein
quoted in the text. The authorities on whom I have drawn most heavily are
listed below.

\paragraph{Non-Euclidean Geometry}
\begin{itemize}
\renewcommand{\labelitemi}{\(\square\)}
  \item R.\ Bonola, \emph{Non-Euclidean Geometry} (Dover Publications, New
  York, 1955).
  \item G.\ Sarton, \emph{Ancient Science and Modern Civilization} (Yale
  University Press, New Haven, 1951), Chapter I.
  \item H.\ Weyl, \emph{Space, Time, Matter}, 4th ed.\ (Dover Publications,
  New York, 1950), Chapter II.
\end{itemize}

\paragraph{Gravitation}
\begin{itemize}
\renewcommand{\labelitemi}{\(\square\)}
  \item F.\ Cajori, historical and explanatory appendix to Isaac Newton's
  \emph{Philosophiae Naturalis Principia Mathematica} (University of
  California Press, 1966).
  \item E.\ Guth, in \emph{Relativity---Proceedings of the Relativity
  Conference in the Midwest}, ed.\ by M.\ Carmeli, S.\ I.\ Fickler, and
  L.\ Witten (Plenum Press, New York, 1970), p.\ 161.
  \item M.\ Jammer, \emph{Concepts of Force} (Harper and Brothers, New York,
  1962), Chapters IV--VII.
  \item E.\ Whittaker, \emph{A History of the Theories of Aether and
  Electricity} (Thomas Nelson and Sons, Edinburgh, 1953), Vol.\ II, Chapter V.
  \item W.\ P.\ D.\ Wightman, \emph{The Growth of Scientific Ideas} (Yale
  University Press, New Haven, 1951), Chapters VIII, X.
\end{itemize}

\paragraph{Relativity}
\begin{itemize}
\renewcommand{\labelitemi}{\(\square\)}
  \item G.\ Holton, ``On the Origins of the Special Theory of Relativity,''
  \emph{Am.\ J.\ Phys.}, \textbf{28}, 627 (1960).
  \item A.\ Koyr\'e, \emph{From the Closed World to the Infinite Universe}
  (Harper and Row, New York, 1958), Chapters VII, IX--XI.
  \item C.\ M{\o}ller, \emph{The Theory of Relativity} (Oxford University
  Press, London, 1952), Chapter I.
  \item W.\ Pauli, \emph{Theory of Relativity} (Pergamon Press, Oxford,
  1958), Parts I, IV.50.
  \item E.\ Whittaker, \emph{A History of Aether and Electricity} (Thomas
  Nelson and Sons, Edinburgh, 1953), Vol.\ I, Chapters VIII--X, XIII; Vol.\
  II, Chapters II, V.
\end{itemize}

\chapterbackmatter{References}
\begin{enumerate}
  \item \phantomsection\label{ch1-ref-1}
  A.\ Einstein, \emph{Annalen der Phys.}, \textbf{49}, 769 (1916). For an
  English translation, see \emph{The Principle of Relativity} (Methuen, 1923,
  reprinted by Dover Publications), p.\ 35.

  \item \phantomsection\label{ch1-ref-2}
  The leading English edition is \emph{Euclid's Elements}, translated with an
  introduction and commentary by T.\ L.\ Heath (rev.\ ed., Cambridge, 1926).

  \item \phantomsection\label{ch1-ref-3}
  These quotations are taken from George Sarton, \emph{Ancient Science and
  Modern Civilization} (University of Nebraska Press, 1954; reprinted by
  Harper and Brothers, New York, 1959), p.\ 26.

  \item \phantomsection\label{ch1-ref-4}
  Quoted by R.\ Bonola, in \emph{Non-Euclidean Geometry}, trans.\ by H.\ S.\
  Carslaw (Dover Press, 1955), pp.\ 65--67.

  \item \phantomsection\label{ch1-ref-5}
  F.\ Klein, \emph{Math.\ Ann.}, \textbf{4}, 573 (1871); \textbf{6}, 112
  (1873); \textbf{37}, 544 (1890); quoted by H.\ Weyl, in
  \emph{Space-Time-Matter}, trans.\ by H.\ L.\ Brose (Dover Press, 1952),
  p.\ 80. A Euclidean model for the Gauss--B\'olyai--Lobachevski geometry was
  given in 1868 by E.\ Beltrami, \emph{Saggio di interpretazione della
  geometria non-euclidea}, quoted by J.\ D.\ North in \emph{The Measure of the
  Universe} (Oxford, 1965), p.\ 60.

  \item \phantomsection\label{ch1-ref-6}
\begingroup
\emergencystretch=1em
  Isaac Newton, \emph{Philosophiae Naturalis Principia Mathematica}, trans.\
  by Andrew Motte, revised and annotated by F.\ Cajori (University of
  California Press, 1966), p.\ 546.
\par
\endgroup

  \item \phantomsection\label{ch1-ref-7}
  R.\ v.\ E\"otv\"os, \emph{Math.\ nat.\ Ber.\ Ungarn}, \textbf{8}, 65
  (1890); R.\ v.\ E\"otv\"os, D.\ Pek\'ar, and E.\ Fekete, \emph{Ann.\
  Phys.}, \textbf{68}, 11 (1922). Also see J.\ Renner, \emph{Hung.\ Acad.\
  Sci.}, Vol.\ 53, Part II (1935).

  \item \phantomsection\label{ch1-ref-8}
  See, for example, A.\ Einstein, \emph{The Meaning of Relativity} (2nd ed.,
  Princeton, 1946), p.\ 56.

  \item \phantomsection\label{ch1-ref-9}
  T.\ D.\ Lee and C.\ N.\ Yang, \emph{Phys.\ Rev.}, \textbf{98}, 1501
  (1955).

  \item \phantomsection\label{ch1-ref-10}
  R.\ H.\ Dicke, in \emph{Relativity, Groups, and Topology}, ed.\ by C.\
  DeWitt and B.\ S.\ DeWitt (Gordon and Breach, New York, 1964), p.\ 167;
  P.\ G.\ Roll, R.\ Krotkov, and R.\ H.\ Dicke, \emph{Ann.\ Phys.\ (N.Y.)},
  \textbf{26}, 442 (1967).

  \item \phantomsection\label{ch1-ref-11}
  J.\ W.\ T.\ Dobbs, J.\ A.\ Harvey, D.\ Paya, and H.\ Horstmann,
  \emph{Phys.\ Rev.}, \textbf{139}, B756 (1965).

  \item \phantomsection\label{ch1-ref-12}
  F.\ C.\ Witteborn and W.\ M.\ Fairbank, \emph{Phys.\ Rev.\ Letters},
  \textbf{19}, 1049 (1967).

  \item \phantomsection\label{ch1-ref-13}
  The most accessible edition is that of Florian Cajori,
  ref.~\ref{ch1-ref-6}.

  \item \phantomsection\label{ch1-ref-14}
  S.\ Newcomb, \emph{Astronomical Papers of the American Ephemeris},
  \textbf{1}, 472 (1882).

  \item \phantomsection\label{ch1-ref-15}
  S.\ Newcomb, article on Mercury in \emph{The Encyclopaedia Britannica},
  11th ed., XVIII, 155 (1910--1911).

  \item \phantomsection\label{ch1-ref-16}
  Ref.~\ref{ch1-ref-6}, p.\ 6 (a different translation is quoted here).

  \item \phantomsection\label{ch1-ref-17}
  \emph{Ibid.}, p.\ 10.

  \item \phantomsection\label{ch1-ref-18}
  G.\ H.\ Alexander, \emph{The Leibniz--Clarke Correspondence} (Manchester
  University Press, 1956). Excerpts are quoted by A.\ Koyr\'e in
  \emph{From the Closed World to the Infinite Universe} (Harper and Row,
  New York, 1958), Chapter XI. (See especially Leibniz's fifth letter.)

  \item \phantomsection\label{ch1-ref-19}
  E.\ Mach, \emph{The Science of Mechanics}, trans.\ by T.\ J.\ McCormack
  (2nd ed., Open Court Publishing Co., 1893).

  \item \phantomsection\label{ch1-ref-20}
  L.\ I.\ Schiff, \emph{Rev.\ Mod.\ Phys.}, \textbf{36}, 510 (1964);
  G.\ M.\ Clemence, \emph{Rev.\ Mod.\ Phys.}, \textbf{19}, 361 (1947);
  \textbf{29}, 2 (1957).

  \item \phantomsection\label{ch1-ref-21}
  James Clark Maxwell, article on ``Ether'' in \emph{The Encyclopaedia
  Britannica}, 9th ed.\ (1875--1889), reprinted in \emph{The Scientific
  Papers of James Clark Maxwell}, ed.\ by W.\ D.\ Niven (Dover Publications,
  1965), p.\ 763. Also see Maxwell's \emph{Treatise on Electricity and
  Magnetism}, Vol.\ II (Dover Publications, 1954), pp.\ 492--493.

  \item \phantomsection\label{ch1-ref-22}
  For an account of these experiments, see C.\ M{\o}ller, \emph{The Theory of
  Relativity} (Oxford Press, London, 1952), Chapter I.

  \item \phantomsection\label{ch1-ref-23}
  A.\ A.\ Michelson and E.\ W.\ Morley, \emph{Am.\ J.\ Sci.}, \textbf{34},
  333 (1887), reprinted in \emph{Relativity Theory: Its Origins and Impact on
  Modern Thought}, ed.\ by L.\ Pearce Williams (John Wiley and Sons, New York,
  1968).

  \item \phantomsection\label{ch1-ref-24}
  T.\ S.\ Jaseja, A.\ Javan, J.\ Murray, and C.\ H.\ Townes, \emph{Phys.\
  Rev.}, \textbf{133}, A1221 (1964).

  \item \phantomsection\label{ch1-ref-25}
  G.\ F.\ Fitzgerald, quoted by O.\ Lodge, \emph{Nature}, \textbf{46}, 165
  (1892). Also see O.\ Lodge, \emph{Phil.\ Trans.\ Roy.\ Soc.},
  \textbf{184A} (1893).

  \item \phantomsection\label{ch1-ref-26}
  H.\ A.\ Lorentz, \emph{Zittingsverslagen der Akad.\ van Wetenschappen},
  \textbf{1}, 74 (November 26, 1892); \emph{Versuch einer Theorie der
  elektrischen und optische Erscheinungen in bewegten K\"orpern} (E.\ J.\
  Brill, Leiden, 1895); \emph{Proc.\ Acad.\ Sci.\ Amsterdam} (English
  version), \textbf{6}, 809 (1904). The third reference and a translated
  excerpt from the second are available in \emph{The Principle of
  Relativity}, ref.~\ref{ch1-ref-1}.

  \item \phantomsection\label{ch1-ref-27}
  J.\ H.\ Poincar\'e, \emph{Rapports pr\'esent\'es au Congr\`es
  International de Physique r\'euni \`a Paris} (Gauthier Villiers, Paris,
  1900); speech at the St.\ Louis International Exposition in 1904, trans.\
  by G.\ B.\ Halstead, \emph{The Monist}, \textbf{15}, 1 (1905), reprinted
  in \emph{Relativity Theory: Its Origins and Impact on Modern Thought},
  ref.~\ref{ch1-ref-23}; \emph{Rend.\ Circ.\ Mat.\ Palermo}, \textbf{21}, 129
  (1906).

  \item \phantomsection\label{ch1-ref-28}
  Sir Edmund Whittaker, \emph{A History of The Theories of Aether and
  Electricity}, Vol.\ II (Thomas Nelson and Sons, London, 1953), Chapter I.

  \item \phantomsection\label{ch1-ref-29}
  For a balanced view of this question, see G.\ Holton, \emph{Am.\ J.\
  Phys.}, \textbf{28}, 627 (1960), reprinted in part in \emph{Relativity
  Theory: Its Origins and Impact on Modern Thought}, ref.~\ref{ch1-ref-23}.

  \item \phantomsection\label{ch1-ref-30}
  A.\ Einstein, \emph{Ann.\ Physik}, \textbf{17}, 891 (1905); \textbf{18},
  639 (1905). Translations are given in \emph{The Principle of Relativity},
  ref.~\ref{ch1-ref-1}.

  \item \phantomsection\label{ch1-ref-31}
  G.\ Holton, ref.~\ref{ch1-ref-29}, and \emph{Isis}, \textbf{60}, 133
  (1969).

  \item \phantomsection\label{ch1-ref-32}
  See ref.~\ref{ch1-ref-30}, and also A.\ Gr\"unbaum, in \emph{Current Issues
  in the Philosophy of Science}, ed.\ by H.\ Feigl and G.\ Maxwell (Holt,
  Rinehart, and Winston, New York, 1961), reprinted in part in
  \emph{Relativity Theory: Its Origins and Impact on Modern Thought},
  ref.~\ref{ch1-ref-23}.

  \item \phantomsection\label{ch1-ref-33}
  A.\ Einstein, \emph{Jahrb.\ Radioakt.}, \textbf{4}, 411 (1907); also see
  M.\ Planck, \emph{Sitzungsber.\ preuss.\ Akad.\ Wiss.}, June 13, 1907,
  p.\ 542; \emph{Ann.\ Phys.\ Leipzig}, \textbf{26} (1908).

  \item \phantomsection\label{ch1-ref-34}
  A.\ Einstein, \emph{Ann.\ Phys.\ Leipzig}, \textbf{35}, 898 (1911). For an
  English translation, see \emph{The Principle of Relativity},
  ref.~\ref{ch1-ref-1}.

  \item \phantomsection\label{ch1-ref-35}
  A.\ Einstein, \emph{Ann.\ Phys.\ Leipzig}, \textbf{38}, 355, 443 (1912).

  \item \phantomsection\label{ch1-ref-36}
  M.\ Abraham, \emph{Lincei Atti}, \textbf{20}, 678 (1911); \emph{Phys.\ Z.},
  \textbf{13}, 1, 4, 176, 310, 311, 793 (1912); \emph{Nuovo Cimento},
  \textbf{4}, 459 (1912).

  \item \phantomsection\label{ch1-ref-37}
  G.\ Nordstr\"om, \emph{Phys.\ Z.}, \textbf{13}, 1126 (1912); \emph{Ann.\
  Phys.\ Leipzig}, \textbf{40}, 856 (1913); \textbf{42}, 533 (1913);
  \textbf{43}, 1101 (1914); \emph{Phys.\ Z.}, \textbf{15}, 375 (1914);
  \emph{Ann.\ Acad.\ Sci.\ fenn.}, \textbf{57} (1914, 1915).

  \item \phantomsection\label{ch1-ref-38}
  A.\ Einstein, \emph{Phys.\ Z.}, \textbf{14}, 1249 (1913); A.\ Einstein and
  M.\ Grossmann, \emph{Z.\ Math.\ Phys.}, \textbf{62}, 225 (1913);
  \textbf{63}, 215 (1914); A.\ Einstein, \emph{Vierteljahr Nat.\ Ges.\
  Z\"urich}, \textbf{58}, 284 (1913); \emph{Archives sci.\ phys.\ nat.},
  \textbf{37}, 5 (1914); \emph{Phys.\ Z.}, \textbf{14}, 1249 (1913).

  \item \phantomsection\label{ch1-ref-39}
  A.\ Einstein, \emph{Sitzungsber.\ preuss.\ Akad.\ Wiss.}, 1914, p.\ 1030;
  1915, pp.\ 778, 799, 831, 844. Also see D.\ Hilbert, \emph{Nachschr.\ Ges.\
  Wiss.\ G\"ottingen}, November 20, 1915, p.\ 395.

  \item \phantomsection\label{ch1-ref-40}
  This famous experiment may in fact be a myth. See A.\ I.\ Miller,
  \emph{Isis}, to be published (1972).
\end{enumerate}
